\scanpage{293}
\mspage{34}
i.e. tous deux directs, ou tous deux rétrogrades.

Ainsi, si on regarde la suite $s_0,s_1,\ldots,s_{\unsure{m|n}}$
des sommets critiques de \unsure{$\tau$|$I$}, alors \deleted{\unsure{$f$ les $f(s_i)$}}
\[
t_0=f(s_0),\quad t_1=f(s_1)\quad\ldots\quad t_m=f(s_m)
\]
forment une suite de sommets de $I_{\unsure{\tau|t}'}$ \deleted{qui} telle
que deux consécutifs soient reliés par une flèche
de $I_{\tau'}$. Donc la sous-catégorie pleine de
$I_{\tau'}$ définie par ces $t_i$ est \unsure{connexe|convexe}, c'est
donc un sous-intervalle. Si $f$ \deleted{\unsure{transforme}}\deleted{\unsure{est un}}\inserted{est}
croissant, ce sous intervalle est $[t_0,t_m]$.
\deleted{morphisme d'inter}
\deleted{En tous les cas} Si $f$ est un morphisme
d'intervalles, il contient tous les sommets
critiques de $I_{\tau'}$. Dans tous les cas, \deleted{\unsure{tout $t_i$}}
\deleted{\unsure{Si, si}}\deleted{sommet critique \unsure{en}}
prenant $t_i,t_j$ $(i<j)$ et appliquant ce qui précède à $f\mid J$, où $J$ est
le sous-intervalle de $I$ d'origine $s_i$, d'extrémité $s_j$, on voit que tout sommet critique
de $I_{\tau'}$ compris entre $t_i$ et $t_j$ est \deleted{\unsure{intégré}}\inserted{\unsure{l'un|un}}
\unsure{des|les} $t_k$ $(i\leq k\leq j)$. Pour résumer :

\uline{Cor 1} Soient $s_i\unsure{<|\leq}s_j$ deux sommets critiques de
\unsure{$I_{\tau'}$|$I_\tau$}. Alors \deleted{\unsure{pour}} tout sommet critique de $I_{\tau'}$
compris entre $f(s_i)=t_i$ et $f(s_j)=t_j$, \unsure{est} \deleted{\unsure{un}}
l'image d'un sommet critique \inserted{$s_k$} compris entre
$s_i$ et $s_j$ (i.e. $i\leq k\leq j$). En particulier,
si $f$ est un morphisme d'intervalles
\[
f(I_{\tau\,\mathrm{red}})\supset I_{\tau'\,\mathrm{red}}
\]
(où $I_{\tau\,\mathrm{red}}$ désigne l'ens des sommets critiques \inserted{de $I_\tau$},
ou aussi le sous-\deleted{catégorie}intervalle réduit qu'il définit).
% archive sheet number: 292
